\documentclass[11pt]{letter}

\usepackage[margin=0.5in]{geometry}
\usepackage{hyperref}

\begin{document}
\begin{letter}{}
\opening{Dear Hiring Team,}

I am writing to apply for the Machine Learning Engineer position, Job ID 10350, at Wawanesa Insurance. I completed my M.Sc. in Computer Science at the University of Calgary and bring experience in Python/PyTorch model development, large-scale data processing, and research communication. I am interested in the opportunity to turn analytical work into useful business decisions, with the care around evaluation and clear explanation that insurance applications require. I am based in Calgary and would welcome working with the team through your hybrid arrangement.

My graduate machine learning work focused on digital elevation model super-resolution. I developed and trained deep ResNet models in PyTorch, built data acquisition and preprocessing workflows, and evaluated reconstruction quality using numerical and terrain-specific measures. A central issue was whether evaluation reflected performance on unfamiliar geography. I identified how random splitting could produce optimistic results and used region-based validation to assess generalization more realistically. This taught me to look beyond a single score, examine the assumptions behind an experiment, and communicate the limits of the evidence alongside the results.

That approach is relevant to Wawanesa's focus on interpretable and fair machine learning. My work has involved understanding where models perform well, where they fail, and how evaluation choices affect the conclusions drawn. I would bring those habits to learning the insurance domain and the team's approach to interpretability and fairness analysis. I am particularly interested in working with analysts and other stakeholders to define what a model needs to accomplish before choosing how to measure success.

My other major graduate project, undertaken with BigGeo and Mitacs, involved designing a framework for managing and retrieving large spatiotemporal Earth observation datasets. I built C++/Python data pipelines, tensor-based storage, and multiresolution retrieval methods. C++ multithreading reduced encoding time for ten tensors from 233 seconds to 26 seconds, while our scaled Canada-wide representation reduced storage by 38\%. This research led to a first-author publication in Remote Sensing. It strengthened my ability to design experiments, compare performance, and explain technical results in a form that others can assess and build on.

I also bring backend and database experience. At Asr Gooyesh Pardaz, I developed Django REST APIs backed by PostgreSQL for NLP and speech services, implemented authentication and payment-integration features, and contributed to a chatbot system. As a teaching assistant at the University of Calgary, I supported students working on SQL, database design, AWS-based data pipelines, and large-scale processing. More recently, I mentored students through Neuromatch Academy's Deep Learning course. These experiences have made explaining technical concepts and helping others work through problems a regular part of my practice.

Wawanesa's combination of analytical work, collaboration, and responsibility for decisions that affect its members appeals to me. I would welcome the opportunity to contribute my research, model-evaluation, and software development experience while developing a deeper understanding of the business problems your team is solving.

Thank you for considering my application. I look forward to discussing how I could contribute to Wawanesa's machine learning work.

\closing{Sincerely,

Amir Mirzai Golpayegani

\href{mailto:amirgolpaa24@gmail.com}{amirgolpaa24@gmail.com}}
\end{letter}
\end{document}
